\documentclass[../../main.tex]{subfiles}
\graphicspath{{../../image/}{../../figure/}}
\begin{document}

\section{松弛}\label{section:ch02-松弛}

前几节引入的连续介质物理状态变量，代表与材料分子相联系的某些物理量（例如速度、平动动能、转动动能、化学能等）的统计平均。这些量演化并最终“松弛”到局部平衡态，后者由能量均分以及统计物理定律所规定的其他条件来表征。当局部平衡达到的时间尺度远小于物体整体运动的时间尺度时，前几节所考察的热弹性本构关系是适用的。相反，当松弛时间与运动时间尺度同阶时，即使在连续介质物理的框架内也必须考虑松弛机制。做法是引入额外的\emph{内部状态变量}，用以度量偏离局部平衡的程度。局部平衡态张成一个嵌入在扩展状态空间中的流形。内部状态变量满足特殊的本构关系，其形式为带耗散源项的平衡律，这些源项驱使状态向量趋向局部平衡。

文献中讨论了种类繁多的松弛理论；读者可通过下面的例子一窥其共同的底层结构。

我们考察一个不导热的连续介质，其等熵响应由本构关系
\begin{align}
\varepsilon=\hat{\varepsilon}(F,\Sigma),
\label{align:ch02-2-8-1}
\end{align}
\begin{align}
S=P(F)+\rho_{0}\Sigma,
\label{align:ch02-2-8-2}
\end{align}
给出，其中 $\Sigma$ 是取值于 $\mathbb{M}^{m\times m}$ 的内部变量，并满足如下形式的平衡律：
\begin{align}
\rho_{0}\dot{\Sigma}=\frac{1}{\tau}[\Pi(\Sigma)-F].
\label{align:ch02-2-8-3}
\end{align}
于是材料表现出瞬时弹性响应（体现在项 $P(F)$ 中），叠加由 $\Sigma$ 的松弛所诱导的粘性响应。正常数 $\tau$ 称为\emph{松弛时间}。

公设：任何平衡了线动量\eqref{align:ch02-2-3-4} 的光滑运动都必须恒等地满足熵不等式\eqref{align:ch02-2-5-32}，由此推出
\begin{align}
S=\rho_{0}\partial_{F}\hat{\varepsilon}(F,\Sigma),
\label{align:ch02-2-8-4}
\end{align}
\begin{align}
\mathrm{tr}\,[\partial_{\Sigma}\hat{\varepsilon}(F,\Sigma)\dot{\Sigma}^{\top}]\leqslant 0.
\label{align:ch02-2-8-5}
\end{align}
将\eqref{align:ch02-2-8-4} 和 \eqref{align:ch02-2-8-5} 与 \eqref{align:ch02-2-8-2}、\eqref{align:ch02-2-8-3} 相结合，我们推出
\begin{align}
\varepsilon=\sigma(F)+\mathrm{tr}\,(\Sigma F^{\top})+h(\Sigma),
\label{align:ch02-2-8-6}
\end{align}
\begin{align}
P(F)=\rho_{0}\partial_{F}\sigma(F),\qquad\Pi(\Sigma)=-\partial_{\Sigma}h(\Sigma).
\label{align:ch02-2-8-7}
\end{align}

当 $h$ 严格凸时，\eqref{align:ch02-2-8-3} 中的源项是耗散性的，并驱使 $\Sigma$ 趋向局部平衡 $\Sigma=H(F)$，其中 $H$ 是 $\Pi$ 的反函数。$\Pi^{-1}$ 存在是因为 $-\Pi$ 严格单调，即
\begin{align}
\mathrm{tr}\,\{[\Pi(\Sigma)-\Pi(\bar{\Sigma})][\Sigma-\bar{\Sigma}]^{\top}\}<0,\qquad\text{对任一 }\Sigma\neq\bar{\Sigma}.
\label{align:ch02-2-8-8}
\end{align}

在局部平衡中，材料响应如同一个具有内能
\begin{align}
\varepsilon=\tilde{\varepsilon}(F)=\sigma(F)+\mathrm{tr}\,[H(F)F^{\top}]+h(H(F))
\label{align:ch02-2-8-9}
\end{align}
和 Piola--Kirchhoff 应力
\begin{align}
S=P(F)+\rho_{0}H(F)=\rho_{0}\partial_{F}\tilde{\varepsilon}(F)
\label{align:ch02-2-8-10}
\end{align}
的弹性材料。

\end{document}
